% =============================================================================
% CONCLUSION AND FUTURE WORK
% =============================================================================
\chapter{Conclusion and Future Work}
\label{ch:conclusion}

\section{Conclusion}

StockSense AI was conceived with a specific and grounded objective: to make inventory
intelligence genuinely accessible to small Indian kirana store owners who currently operate
without any digital support for stock management. The project was not an attempt to build
a complex, heavyweight enterprise system. It was, instead, a deliberate exercise in designing
the simplest possible system that could deliver meaningful, automated, AI-augmented inventory
management within the real financial and operational constraints of the target user segment.

The system successfully achieves all six objectives defined at the outset of the project.
The automated intelligence pipeline — comprising ten modular n8n workflows — calculates
stockout dates for every product using daily sales velocity, classifies products into HIGH,
MEDIUM, and LOW risk tiers, generates fresh stock alerts for at-risk products, invokes the
Groq Llama~3.1~8B large language model to produce natural-language reorder suggestions with
supplier selection reasoning, scores all suppliers using the Weighted Sum Model, and delivers
formatted WhatsApp alerts to the store owner through the Twilio API. The entire pipeline
executes from end to end in approximately thirteen seconds and runs autonomously on a daily
cron schedule without any manual intervention.

The AI vision capabilities — product onboarding from a photograph using Google Gemini~2.5
Flash and supplier invoice scanning with fuzzy inventory matching — represent a significant
reduction in data entry burden for the store owner. Features that would previously have
required manual typing of product names and quantities from physical bills are now automated
through a simple photograph upload.

The React~19 frontend provides a responsive, real-time dashboard with Supabase WebSocket
subscriptions ensuring that pipeline outputs are immediately reflected on screen without any
page reload. The twelve-page application covers every functional area from onboarding to
supplier analytics, presenting complex inventory intelligence in a form that a non-technical
store owner can understand and act upon.

The prototype operates at zero infrastructure cost using free-tier services, and the estimated
production deployment cost of below \rupee{}500 per month validates the economic viability
argument at the core of the project's motivation. Comparative analysis against existing tools
such as Vyapar, Zoho Inventory, and Khatabook confirms that no commercially available solution
at this price point offers the combination of automated stockout prediction, LLM-driven reorder
recommendations, supplier scoring, and WhatsApp alerting that StockSense AI delivers.

\textbf{Limitations.} The system has several acknowledged limitations that should be stated
honestly. First, the stockout prediction formula relies on the \texttt{avg\_daily\_sales} field,
which is currently a manually entered historical average. The system does not compute this
value dynamically from transaction history, meaning it is only as accurate as the value the
store owner provides. Second, the Weighted Sum Model for supplier scoring is deterministic
and static; it does not learn from historical supplier performance feedback or adapt its
weights over time. Third, the demand forecast does not account for seasonal spikes or
festival-driven demand surges, which are significant in the Indian retail context. Fourth,
the current implementation enforces single-store tenancy; the schema is designed for
multi-tenancy but the row-level security policies and tenant-aware API routing required for
production multi-store operation have not been implemented. Fifth, the system has been tested
only against a synthesised dataset of ten products; performance and accuracy at production
scale with hundreds of products and multiple supplier categories remains to be validated.

Despite these limitations, StockSense AI constitutes a functional and deployable proof of
concept that demonstrates the viability of the proposed approach. It shows that the combination
of n8n workflow automation, Supabase real-time databases, Groq LLM inference, and Google
Gemini multimodal AI can be integrated into a cohesive inventory intelligence system at a
cost and complexity level appropriate for small Indian retailers.

\section{Future Work}

The following enhancements are identified as the most impactful directions for future
development of StockSense AI, ordered by estimated implementation effort and user impact:

\textbf{Dynamic Sales Velocity Computation.} The most significant immediate improvement
would be to compute \texttt{avg\_daily\_sales} automatically from the Stock Transactions
audit log rather than relying on a manually entered value. This would make the stockout
prediction formula self-updating and progressively more accurate as transaction history
accumulates. A rolling 30-day average window would be a natural starting point.

\textbf{Seasonal Demand Multiplier.} The Products table already includes \texttt{is\_seasonal}
and \texttt{peak\_season\_months} columns. A future enhancement to WF-02 would apply a
configurable seasonal multiplier to \texttt{avg\_daily\_sales} during identified peak months,
improving stockout prediction accuracy around Diwali, Holi, harvest season, and other
demand spikes that are predictable in the Indian retail calendar.

\textbf{Multi-Store Tenancy with Row-Level Security.} The database schema is already
designed for multi-tenancy through the \texttt{store\_id} foreign key on all data tables.
Enabling Supabase Row Level Security (RLS) policies that filter all queries by the
authenticated user's \texttt{store\_id} would allow the system to serve multiple independent
store owners on a single infrastructure instance, dramatically improving the unit economics
of the platform.

\textbf{Native Mobile Application.} Since most kirana store owners in India interact
with digital services primarily through smartphones, a React Native mobile application
built on the same Supabase backend and n8n API would significantly broaden the system's
accessibility and daily usage frequency.

\textbf{Two-Way WhatsApp Interaction.} The current WhatsApp integration is one-directional:
the system sends alerts to the store owner. A future enhancement using Twilio's inbound
message handling capability would allow the store owner to reply to an alert with a command
(e.g., ``reorder 50 units'') and have the system automatically approve the corresponding
reorder suggestion and log the action, without requiring the owner to open the web application.

\textbf{Automated Purchase Order Dispatch.} The logical extension of an approved reorder
suggestion is the actual placement of an order with the supplier. This could be implemented
by integrating with supplier communication channels (WhatsApp Business API or email) to
send formatted purchase order messages, and potentially with UPI payment APIs to facilitate
automated payment initiation for known supplier accounts.

\textbf{Adaptive Supplier Scoring.} Replacing the static Weighted Sum Model weights with
a learning system that adjusts the importance of reliability, price, and delivery speed
based on the store owner's historical approval patterns for reorder suggestions would make
the supplier scoring more personalised and contextually relevant over time.

These future directions collectively outline a path from the current prototype to a
production-grade, multi-tenant SaaS platform — StockSense AI as it could eventually serve
not just individual kirana stores but the broader ecosystem of small Indian retailers.
